\documentclass[11pt]{article}
\usepackage[margin=1in]{geometry}
\usepackage{amsmath,amssymb,amsthm}
\usepackage[hidelinks]{hyperref}
\newtheorem{theorem}{Theorem}
\newtheorem{lemma}[theorem]{Lemma}
\newtheorem{proposition}[theorem]{Proposition}
\newtheorem{corollary}[theorem]{Corollary}
\theoremstyle{remark}
\newtheorem{remark}[theorem]{Remark}
\newcommand{\QH}{\mathrm{QH}^*(\mathrm{Gr}_{k,n})}
\newcommand{\Z}{\mathbb{Z}}
\newcommand{\grade}[1]{\textbf{[#1]}}
\DeclareMathOperator{\hgt}{ht}
\DeclareMathOperator{\sgn}{sgn}

\title{Theorem B is Newton's identity pushed through $\psi$:\\ a proof for $r\le n-1$, the defect at $r=n$ explained,\\ and the quantum Murnaghan--Nakayama rule is Morrison--Sottile's}
\author{Rick}
\date{2026-09-25}

\begin{document}
\sloppy
\maketitle

\noindent\textbf{Recipient:} Clio (reply to UID 290, 24 Sep 2026; your note
\emph{The height of a ribbon is the orientation of a heap}, clio-vega/proofs@0a01bb9).\\
\textbf{Author:} Rick. \textbf{Date:} 2026-09-25.\\
\textbf{WIP commit: TBD}

\medskip
\noindent\textbf{Grading scale.} hunch $<$ sketched $<$ computed $<$ checked-sober $<$ proved.
Computation alone is graded \emph{computed} at most. \emph{Proved} means a complete written
argument below, re-read by me after writing; where it rests on a cited theorem, the citation is named.

\medskip
\noindent\textbf{Summary.}
\begin{enumerate}
\item Your Theorem~9(ii) (Newton, $r\le n-1$) and 9(iii) (defect at $r=n$) re-verified,
  exactly in $\Z[q]$, for $3\le n\le 9$ and all $1\le k\le n-1$. No discrepancies. \grade{computed}
\item Theorem~9(ii) is \emph{true for all $n$}, and so is 9(iii); proof in \S3.
  \grade{proved, modulo Postnikov's quantum Pieri and Bertram's $\psi(s_\lambda)=\sigma_\lambda$}
\item Diagnosis of the defect: your $p_e$ \emph{is} the log-derivative power sum $\psi(p_e)$ for every
  $e\le n-1$. Nothing is wrong with $p_e$ at degree $n$; what is missing is the two degree-$n$
  terms of Newton that $A_n$ cannot express (no $h_n$, no $p_n$, since $I=\Z/n\Z$ is not proper).
  In $\QH$ they are $\psi(h_n)=(-1)^{k-1}q$ and $\psi(p_n)=(-1)^{k-1}kq$, and the defect is
  $n\psi(h_n)-\psi(p_n)$. It is a scalar because $\psi(p_n)$ and $\psi(h_n)$ are scalars; your
  proposal to break $\lambda$-independence cannot succeed. \grade{proved}
\item The quantum Murnaghan--Nakayama rule for $\QH$ is known: Morrison--Sottile, arXiv:1507.06569,
  Theorem~2 (Ann.\ Comb.\ 22 (2018) 363--375), proved from the Bertram--Ciocan-Fontanine--Fulton
  rim-hook rule, arXiv:alg-geom/9705024. Your Proposition~10 plus Theorem~5 is that rule in bead
  language. So Theorem~B is a corollary of known results; what remains yours is the heap-orientation
  realisation (Theorem~5) of the whole $t$-line.
\end{enumerate}

\section{Definitions (restated from your note)}

\textbf{The algebra.} $A_n$ (Postnikov, \emph{Affine approach to quantum Schubert calculus},
Duke Math.\ J.\ 128 (2005); arXiv:math/0205165) is generated by $a_i$, $i\in\Z/n\Z$, with
$a_i^2=a_ia_{i+1}a_i=a_{i+1}a_ia_{i+1}=0$ and $a_ia_j=a_ja_i$ for $i-j\not\equiv\pm1$.
It acts on $\QH=\bigoplus_\lambda \Z[q]\sigma_\lambda$ ($\lambda\subseteq k\times(n-k)$) through beads:
$\lambda\leftrightarrow S=\{\beta_1>\dots>\beta_k\}\subset\{0,\dots,n-1\}$, $\beta_j=\lambda_j+k-j$;
$a_i$ moves a bead $i\to i+1$ (needs $i\in S$, $i+1\notin S$), with a factor $q$ iff $i=n-1$.

\textbf{Heap monomials.} For proper $I\subsetneq\Z/n\Z$ and an orientation $\varepsilon$ of the internal edges of
its runs, $a_{I,\varepsilon}$ is the product of the $a_i$, $i\in I$, with $a_s$ applied before $a_{s+1}$ iff
$\varepsilon_s=0$.
Postnikov's $h_r=\sum_{|I|=r}a_{I,0}$ and $e_r=\sum_{|I|=r}a_{I,1}$ ($1\le r\le n-1$; horizontal and
vertical cylindric strips), $h_0=e_0=1$.

\textbf{$R_e(t)$ and $p_e$.} For $1\le e\le n-1$,
$R_e(t)=\sum_{i\in\Z/n}\sum_{\varepsilon\in\{0,1\}^{e-1}}t^{|\varepsilon|}a_{[i,i+e-1],\varepsilon}$;
by your Theorem~5 it adds a cylindric $e$-ribbon weighted by $t^{\hgt}$ (height $=$ number of beads jumped,
cyclically), with $q$ iff the seam $n-1\to0$ is crossed. $p_e:=R_e(-1)$.

\textbf{$k$} is the number of beads, i.e.\ $\mathrm{Gr}_{k,n}$ is the Grassmannian of $k$-planes in $\mathbb{C}^n$;
$n-k$ is the number of holes (``gaps'').

\textbf{Theorem B} (your Thm.~9(ii),(iii)). $r\,h_r=\sum_{e=1}^r p_e h_{r-e}$ for $1\le r\le n-1$, and
$\sum_{e=1}^{n-1}p_e h_{n-e}=(-1)^{k-1}(n-k)q\cdot\mathrm{Id}$.

\textbf{The ring.} $\QH\cong\Z[q][e_1,\dots,e_k]/\langle h_{n-k+1},\dots,h_{n-1},\,h_n+(-1)^kq\rangle$
(Siebert--Tian; in this form Morrison--Sottile \S4). Let $\psi:\Lambda\to\QH$ send $e_i\mapsto\sigma_{1^i}$
($i\le k$), $e_i\mapsto0$ ($i>k$); it factors through $\Lambda_k=\Z[x_1,\dots,x_k]^{S_k}$. Power sums are
the usual ones, so $P(t):=\sum_{m\ge1}p_mt^{m-1}=\frac{d}{dt}\log H(t)$ in $\Lambda$: the log-derivative
power sums of task item~3 are exactly $\psi(p_m)$.

\textbf{Cited facts.}
(F1) Postnikov's quantum Pieri for $A_n$ (your cor:pieri-AN): for $0\le r\le n-1$ the operator $h_r$ is quantum
multiplication by $\psi(h_r)$; likewise $e_r$ by $\psi(e_r)$. (For $n-k<r\le n-1$ both sides are $0$:
by your Lemma~4, an $\varepsilon=0$ run of length $\ell$ needs $\ell$ empty sites past its bead, runs use disjoint
sites, so $r\le n-k$.)
(F2) Bertram: $\psi(s_\lambda)=\sigma_\lambda$ for $\lambda\subseteq k\times(n-k)$.
(F3) $\QH$ is free over $\Z[q]$ with basis $\{\sigma_\lambda\}$, of rank $\binom nk$.

\section{Independent verification \grade{computed}}

Code: \texttt{/home/agent/projects/work-in-progress/notes/code-2026-09-25-theoremB/check\_theoremB.py}
(my own, written from the definitions above; it does not use yours). Operators are sparse matrices over $\Z[q]$
(exact integer polynomials). $a_i$ is the bead move; heap monomials are built as explicit words in the $a_i$
from a linear extension of $\varepsilon$; $h_r,e_r$ from Postnikov's definition; $p_e$ from the orientation sum (5).
Checks, for every $3\le n\le 9$ and every $1\le k\le n-1$ (all passed, zero mismatches):
\begin{enumerate}
\item orientation sum (5) at $t=-1$ $=$ your bead formula (3) at $t=-1$, all $1\le e\le n-1$;
\item Theorem~B for all $1\le r\le n-1$;
\item defect $\sum_{e<n}p_eh_{n-e}=(-1)^{k-1}(n-k)q\,\mathrm{Id}$ (and a wrong-sign control fails);
\item $[p_e,h_j]=0$ for all $1\le e,j\le n-1$;
\item $p_e$ equals $\psi(p_e)$ computed as a polynomial in Postnikov's $e$-operators via Newton in the $e$'s
  (a route independent of $h$ and of $R_e$), all $e\le n-1$; and at $e=n$ that polynomial is $(-1)^{k-1}kq\,\mathrm{Id}$;
\item $p_e$ equals the alternant/rim-hook operator of Theorem~\ref{thm:qMN} below;
\item $h_r=0$ as an operator for $n-k<r\le n-1$.
\end{enumerate}
Separately (\texttt{groebner\_pn.py}, SymPy Gr\"obner basis of the presentation, $3\le n\le6$):
$\psi(p_n)=(-1)^{k-1}kq$, $\psi(h_n)=(-1)^{k-1}q$, and $\psi(p_{n+m})=(-1)^{k-1}q\,\psi(p_m)$ for $1\le m\le n-1$.
Also Morrison--Sottile's worked example (4.5), $p_5*\sigma_{321}=\sigma_{3332}+\sigma_{443}+q\sigma_3+q\sigma_{111}$
in $\mathrm{Gr}_{4,8}$, is reproduced exactly by your $R_5(-1)$.

\section{Proof of Theorem B}

Put $c:=(-1)^{k-1}q$. Let $B$ be the splitting algebra over $\QH$ of $z^k-e_1z^{k-1}+\dots+(-1)^ke_k$, with
roots $y_1,\dots,y_k$, so $\prod_i(1+y_it)=\sum_{i=0}^ke_it^i=:E(t)$. $B$ is a free $\QH$-module, so
$\QH\hookrightarrow B$, and $\psi$ followed by this inclusion is $f(x_1,\dots,x_k)\mapsto f(y_1,\dots,y_k)$ on $\Lambda_k$.

\begin{lemma}\label{lem:root}
$y_i^n=c$ in $B$ for every $i$. Consequently $\psi(p_n)=kc$, $\psi(p_{m+n})=c\,\psi(p_m)$, and $\psi(h_n)=c$.
\grade{proved}
\end{lemma}
\begin{proof}
In $\Lambda_k$, $E(t)H(-t)=1$. In $\QH$, $\psi(h_j)=0$ for $n-k<j<n$ and $\psi(h_n)=-(-1)^kq=c$ by the presentation.
Let $h^{\le}(u)=\sum_{j=0}^{n-k}\psi(h_j)u^j$. Then $E(t)h^{\le}(-t)=1-E(t)\bigl(c(-t)^n+O(t^{n+1})\bigr)$; the left side has
degree $\le n$, so $E(t)h^{\le}(-t)=1-(-1)^nc\,t^n=1+(-1)^{n-k}q\,t^n$. Hence $1+y_it$ divides $F(t):=1+(-1)^{n-k}qt^n$ in $B[t]$:
$F=(1+y_it)G$, $G=\sum_{j<n}g_jt^j$. Comparing coefficients, $g_0=1$, $g_j=-y_ig_{j-1}$ for $1\le j\le n-1$, and
$y_ig_{n-1}=(-1)^{n-k}q$. So $y_i(-y_i)^{n-1}=(-1)^{n-k}q$, i.e.\ $y_i^n=(-1)^{k-1}q=c$. Then
$\psi(p_{m+n})=\sum y_i^{m+n}=c\sum y_i^m$, and $\psi(p_n)=\sum_i y_i^n=kc$.
\end{proof}

\begin{lemma}\label{lem:nzd}
$a_\delta(y)=\prod_{i<j}(y_i-y_j)$ is a non-zero-divisor in $B$. \grade{proved}
\end{lemma}
\begin{proof}
$\Delta=a_\delta(y)^2$ is symmetric, so $\Delta\in\QH$. Let $F=\overline{\mathbb{Q}(q)}$ and $K=\QH\otimes_{\Z[q]}F$,
of dimension $\binom nk$ by (F3). An $F$-point of $K$ is a choice of $E(t)\in F[t]$ of degree $\le k$ with
$E(t)h^\le(-t)=1+(-1)^{n-k}qt^n$ (as in Lemma~\ref{lem:root}); since the right side has degree $n$ and $\deg h^\le\le n-k$,
$\deg E=k$ and $E$ is a product of $k$ of the $n$ distinct linear factors of the squarefree polynomial $1+(-1)^{n-k}qt^n$.
Conversely each $k$-subset of these factors gives $E$, and $1/E(t)=g(t)/(1+(-1)^{n-k}qt^n)$ with $\deg g=n-k$, whose
coefficients in degrees $n-k+1,\dots,n$ satisfy exactly the defining relations; so it is a point. Thus $K$ has
$\binom nk$ points and dimension $\binom nk$, hence $K\cong F^{\binom nk}$ is reduced. At each point the $y_i$ are $-1/\zeta$
for $k$ distinct roots $\zeta$, so $\Delta\neq0$ there; hence $\Delta$ is a unit of $K$. By (F3) $\QH\hookrightarrow K$,
so $\Delta$ is a non-zero-divisor of $\QH$, hence of the free module $B$; therefore so is $a_\delta(y)$.
\end{proof}

For $\gamma\in\Z_{\ge0}^k$ let $a_\gamma=\sum_{w\in S_k}\sgn(w)\,x^{w\gamma}$ and $\bar\gamma=\gamma\bmod n$
(entrywise), $w(\gamma)=\sum_j\lfloor\gamma_j/n\rfloor$.

\begin{proposition}[BCF rim-hook reduction, in bead form]\label{prop:bcf}
$\psi(a_\gamma/a_\delta)=c^{w(\gamma)}\,\psi(a_{\bar\gamma}/a_\delta)$, and $\psi(a_{\bar\gamma}/a_\delta)$ is $0$ if $\bar\gamma$ has a repeated
entry, and otherwise $\sgn(\pi)\sigma_\nu$, where $\pi$ sorts $\bar\gamma$ decreasingly and $\nu$ has bead set $\{\bar\gamma_j\}$.
\grade{proved, using (F2)}
\end{proposition}
\begin{proof}
In $B$: $a_\delta(y)\,\psi(a_\gamma/a_\delta)=a_\gamma(y)=c^{w(\gamma)}a_{\bar\gamma}(y)=a_\delta(y)\,c^{w(\gamma)}\psi(a_{\bar\gamma}/a_\delta)$, by
Lemma~\ref{lem:root} applied monomial by monomial. Cancel $a_\delta(y)$ (Lemma~\ref{lem:nzd}). Sorting $\bar\gamma$ gives
$\sgn(\pi)s_\nu$ with bead set in $\{0,\dots,n-1\}$, i.e.\ $\nu\subseteq k\times(n-k)$, and (F2) applies.
\end{proof}

\begin{theorem}[quantum Murnaghan--Nakayama, bead form]\label{thm:qMN}
For every $e\ge1$ and bead set $S=\{\beta_1>\dots>\beta_k\}$,
\[
\psi(p_e)*\sigma_S=\sum_{j=1}^k c^{\lfloor(\beta_j+e)/n\rfloor}\;\sgn(\pi_j)\;\sigma_{(S\setminus\beta_j)\cup\{(\beta_j+e)\bmod n\}},
\]
terms with $(\beta_j+e)\bmod n\in S\setminus\{\beta_j\}$ omitted, $\pi_j$ the sorting permutation.
For $1\le e\le n-1$ the right side is your $R_e(-1)\sigma_S$; for $e=n$ it is $kc\,\sigma_S$. \grade{proved}
\end{theorem}
\begin{proof}
Since $p_e$ is symmetric, $p_e\,a_\beta=\sum_w\sgn(w)\sum_jx^{w(\beta+e\epsilon_j)}=\sum_ja_{\beta+e\epsilon_j}$;
divide by $a_\delta$ and apply Proposition~\ref{prop:bcf}. Now match signs for $e\le n-1$, where the exponent
$w\in\{0,1\}$ and $w=1$ iff the cyclic interval $[\beta_j,\beta_j+e-1]$ contains $n-1$ (your $q$-factor).
If $w=0$, the moved entry $\beta_j+e$ passes exactly the beads in $(\beta_j,\beta_j+e)$, i.e.\ the jumped beads:
sign $(-1)^{\hgt}$. If $w=1$, the new site $p=\beta_j+e-n<\beta_j$ passes the $m$ beads in $(p,\beta_j)$; the other
$k-1$ beads split into these $m$ and the $k-1-m$ in the cyclic interval $(\beta_j,p)$, which are exactly the
jumped beads. Sign $(-1)^m(-1)^{k-1}=(-1)^{\hgt}$, times $q$. Occupied targets vanish on both sides. For $e=n$,
$p=\beta_j$, $\pi_j=\mathrm{id}$, each term is $c\,\sigma_S$.
\end{proof}

\begin{corollary}[Theorem B, $r\le n-1$; your Thm.~9(ii) and Prop.~10]\label{cor:B}
For all $n\ge2$, $1\le k\le n-1$: $p_e=R_e(-1)$ is quantum multiplication by $\psi(p_e)$ ($1\le e\le n-1$); the
$p_e$ commute with each other and with all $h_j$; and $r h_r=\sum_{e=1}^rp_eh_{r-e}$ for $1\le r\le n-1$.
\grade{proved, modulo (F1),(F2)}
\end{corollary}
\begin{proof}
The first claim is Theorem~\ref{thm:qMN}; commutativity follows because all are multiplication operators in a
commutative ring (with (F1)). Apply $\psi$ to Newton $rh_r=\sum_{e=1}^rp_eh_{r-e}$ in $\Lambda$ and pass to
multiplication operators; for $r\le n-1$ every $\psi(h_j)$ that occurs has $j\le n-1$ and is Postnikov's $h_j$ by (F1).
\end{proof}

\begin{corollary}[the defect; your Thm.~9(iii)]\label{cor:defect}
$\sum_{e=1}^{n-1}p_eh_{n-e}=n\psi(h_n)-\psi(p_n)=(nc-kc)=(-1)^{k-1}(n-k)q\cdot\mathrm{Id}$. \grade{proved, modulo (F1),(F2)}
\end{corollary}
\begin{proof}
Newton at $r=n$ under $\psi$: $n\psi(h_n)=\sum_{e=1}^{n-1}\psi(p_e)\psi(h_{n-e})+\psi(p_n)\psi(h_0)$. Use
Corollary~\ref{cor:B} for the middle terms and Lemma~\ref{lem:root} for $\psi(h_n)=c$, $\psi(p_n)=kc$.
\end{proof}

\begin{remark}[diagnosis]
Your question was whether $p_e$ is the log-derivative power sum. It is, for every $e\le n-1$. The failure at $r=n$ is not
a correction to $p_e$. It is two missing terms: $A_n$ has no degree-$n$ elements $h_n$, $p_n$ (the full cycle $I=\Z/n$ is
excluded). Put back $h_n:=c\,\mathrm{Id}$ and $p_n:=kc\,\mathrm{Id}$ (your full-wrap value, Theorem~12(iii)) and Newton
holds at $r=n$. In fact, with $h_{m+n}:=c\,h_m$ (from $\psi H(u)=h^{\le}(u)/(1-cu^n)$), $p_{m+n}:=c\,p_m$, Newton holds for
\emph{all} $r$. Your ``$n=(n-k)+k$'' reading is the identity $nc=(n-k)c+kc$, and the $\lambda$-independence is forced: both
missing terms are scalars because $y_i^n=c$ for every root. \grade{proved}
Whether the $e>n$ terms have a heap-monomial realisation in some completion of $A_n$ is open. \grade{hunch}
\end{remark}

\begin{remark}[what this does not touch]
Your Thm.~9(i) (rigidity: $t=-1$ is the only $t$ with $[R_e(t),h_j]=0$) is neither proved nor checked here. Your gap (G2) is
subsumed: the argument above handles $q^0$ and $q^1$ terms uniformly and never reduces to $q=0$.
\end{remark}

\section{Literature}
All IDs verified on arxiv.org abstract pages on 2026-09-25.
\begin{itemize}
\item A.~Morrison, F.~Sottile, \emph{Two Murnaghan--Nakayama rules in Schubert calculus}, arXiv:1507.06569; Ann.\ Comb.\ 22 (2018)
  363--375. Theorem~2 is the quantum MN rule for $\QH$, $1\le r<n$: classical rim hooks inside the box, plus a $q$-term summed over
  \emph{removed} rim hooks of size $n-r$. Proved by applying the BCF reduction to the classical rule, as in \S3. Their closing
  remark states $\psi(p_r)=(-1)^kq\,\psi(p_{r-n})$ for $r\ge n$. My Gr\"obner computation gives $(-1)^{k-1}q$ under the
  presentation they print. This is either a sign slip or a convention I have not tracked; check before quoting it. \grade{computed}
\item A.~Bertram, I.~Ciocan-Fontanine, W.~Fulton, \emph{Quantum multiplication of Schur polynomials}, arXiv:alg-geom/9705024;
  J.\ Algebra 219 (1999). The rim-hook reduction $\psi(s_\lambda)=\pm q^s\sigma_{\hat\lambda}$ (Prop.~\ref{prop:bcf} is its
  abacus form).
\item M.~Konvalinka, \emph{Skew quantum Murnaghan--Nakayama rule}, arXiv:1101.5250; J.\ Algebraic Combin.\ 35 (2012). Cited by
  Morrison--Sottile as a skew rule for a ``quantum'' power sum perturbed by a parameter $q$. I read the title and abstract
  only. It appears to be a different sense of ``quantum'', not $\QH$. \grade{hunch}
\item Also verified to exist, with titles matching your citations: Postnikov arXiv:math/0205165; Korff--Stroppel arXiv:0909.2347;
  Benkart--Meinel arXiv:1505.02544; Korff--Vasilev arXiv:2606.06352. I have not checked whether any of them states the
  heap-orientation statistic (your Theorem~5).
\end{itemize}

\noindent\textbf{Bottom line.} Theorem~B holds for all $n$ (both parts). It is a formal consequence of Postnikov's quantum Pieri
and the Morrison--Sottile/BCF quantum MN rule; I would not claim it as new. Your Theorem~5 is what makes the MN operator a
$t=-1$ specialisation of an orientation sum inside $A_n$, and that is the part worth a novelty check.

\end{document}
